\documentclass[a5paper,10pt]{article}

% https://github.com/InDevRus/filippov-solutions

% Нумерация страниц
\pagenumbering{gobble}

% Настройка полей
\usepackage{geometry}
\geometry{left=1cm}
\geometry{right=1cm}
\geometry{top=1cm}
\geometry{bottom=1.5cm}

% Вставка таблицы
\usepackage{array}
\usepackage{tabularx}

% Инструменты выравнивания формул
\usepackage{amsmath}
\usepackage{mathtools}

% Диакритика в математике
\usepackage{amsfonts}

% Вычисления размеров на ходу
\usepackage{calc}

% Удобные записи производных
\usepackage[italic=true]{derivative}

% Настройки сносок
\usepackage{enotez}

% Длинные знаки векторов
\usepackage{anyfontsize}
\usepackage[b]{esvect}

% Вставка картинок
\usepackage{graphicx}

% Вставка красной строки
\usepackage{indentfirst}
\setlength{\parindent}{1em}

% Условия в командах
\usepackage{xifthen}

% Луа-скрипты
\usepackage{luacode}

% Настройка команд отдельно в математическом и текстовом режимах
\usepackage{mathcommand}

% Для повторения LaTeX кода
\usepackage{multido}

% Конфигурация отступов формул
\delimitershortfall-1sp
\usepackage{mleftright}
\mleftright

% Растягивание объектов
\usepackage{relsize}

% Обтекание текстом
\usepackage{wrapfig}
\usepackage{subcaption}
\usepackage{floatflt}

% Поддержка ввода кириллицы
\usepackage[english, russian]{babel}

% Настройка корневой позиции для компилятора
\usepackage{import}
\providecommand*{\ProjectRootPath}{..}

\import{\ProjectRootPath/preambles/macros/}{theorems}

\import{\ProjectRootPath/preambles/fonts}{font_setup}

% Знак градуса
\usepackage{siunitx}

\import{\ProjectRootPath/preambles/macros/}{operators}
\import{\ProjectRootPath/preambles/macros/}{redefinitions}

\import{\ProjectRootPath/preambles/fonts}{letters_setup}
\import{\ProjectRootPath/preambles/fonts/adjustments}{custom_superscripts}

\import{\ProjectRootPath/preambles/custom}{formatting}
\import{\ProjectRootPath/preambles/custom}{frames}
\import{\ProjectRootPath/preambles/custom}{list_setup}

% TODO: перевести команды на современный стиль объявления

\begin{document}
    $xy\` - y = \br{x + y} \ln\br{\dfrac{x + y}{x}}$.
    $y\` = \dfrac{y}{x} + \br{1 + \dfrac{y}{x}} \ln\br{1 + \dfrac{y}{x}}$. \\
    Заменяем
    $y\br{x} = t\br{x} \cdot x$.
    $y\` = t\`x + t$. \\
    $t\`x + t = t + \br{1 + t} \ln\br{1 + t}$.
    $t\`x = \br{1 + t} \ln\br{1 + t}$.
    $\dfrac{t\`}{\br{1 + t}\ln\br{1 + t}} = \dfrac {1}{x}$,
    причём $t\br{x} = 0$ и $t\br{x} = -1$ являются решениями,
    которые дают $y\br{x} = 0$ и $y\br{x} = -x$ соответственно.
    $$\tint \dfrac{\di t}{\br{1 + t}\ln\br{1 + t}}
    = \tint \dfrac{\di\br{\ln\br{1 + t}}}{\ln\br{1 + t}}
    = \ln{\vbr{\ln\br{1 + t}}} + C. $$

    $\ln{\vbr{\ln\br{1 + t}}} = \ln\vbr{x} + C$.
    $\ln\br{1 + t} = C x$.
    $t\br{x} = \pow{e}{Cx} - 1$. \\
    Итак, общее решение равно
    $y\br{x} = x\br{\pow{e}{Cx} - 1}$, а также дополнительно
    $y\br{x} = -x$.
\end{document}
